\section{Dynamic Semantics}
\label{sec:dynamic}
The dynamic semantics of $\minilang$ define a transition
system for evaluating $\minilang$ expressions.
In order to know when we are done evaluating an expression
to a \emph{value}, we need to define values.
Values in $\minilang$ are either a single number or a
single string.
\[
\cinfer{num[$n$] value}{ }
\qquad
\cinfer{str[$s$] value}{ }
\]


The inductive definition of the transition relation for
evaluating expressions is shown in Figure~\ref{fig:eval_rules}.
Rules \code{D.1-3} define how to evaluate additions.
Rule \code{D.3} says that the addition of two single
numbers \emph{steps to} (evaluation step) a number that is the
sum of those two numbers.
Rule \code{D.1} says if an evaluation step can be performed
on the left summand, then the addition expression steps to
an expression that is the same except with the step performed
on the left summand.
Once we are done evaluating the left subexpression to a number,
rule \code{D.2} allows us to evaluate the right subexpression.
Rules \code{D.4-6} are analogous to rules \code{D.1-3}
for string concatenation.
Rule \code{D.7} says we evaluate the expression ($e_1$) that
will be bound to the variable of the \code{let} expression.
Once $e_1$ becomes a value, then rule \code{D.8} says that
we can evaluate the body ($e_2$) of the \code{let} expression
by replacing the variable with that value.

\begin{figure}[htb]
\[
\cinfer[D.1]{+($e_1$; $e_2$) $\stepto$ +($e'_1$; $e_2$)}{e_1 \stepto e'_1}
\qquad
\cinfer[D.2]{+(num[$n_1$]; $e_2$) $\stepto$ +(num[$n_1$]; $e'_2$)}{e_2 \stepto e'_2}
\]

\[
\cinfer[D.3]{+(num[$n_1$]; num[$n_2$]) $\stepto$ num[$n_1 + n_2$]}{ }
\]

\[
\cinfer[D.4]{\^{}($e_1$; $e_2$) $\stepto$ \^{}($e'_1$; $e_2$)}{e_1 \stepto e'_1}
\qquad
\cinfer[D.5]{\^{}(str[$s_1$]; $e_2$) $\stepto$ \^{}(str[$s_1$]; $e'_2$)}{e_2 \stepto e'_2}
\]

\[
\cinfer[D.6]{\^{}(str[$s_1$]; str[$s_2$]) $\stepto$ str[$s_1 \texttt{\^{}} s_2$]}{ }
\]

\[
\cinfer[D.7]{let($x$; $e_1$; $e_2$) $\stepto$ let($x$; $e'_1$; $e_2$)}{e_1 \stepto e'_1}
\qquad
\cinfer[D.8]{let($x$; $e_1$; $e_2$) $\stepto$ $[e_1 / x]e_2$}{e_1 \; \texttt{value}}
\]
\caption{Dynamic Semantics of $\minilang$}
\label{fig:eval_rules} 
\end{figure}


